\documentclass{article}
\usepackage{amsmath, amssymb, amsthm}
\usepackage{hyperref}
\usepackage{geometry}
\geometry{margin=1in}

\title{Erd\H{o}s Problem \#761}
\date{Last updated: 2025-08-31}
\author{Source: erdosproblems.com}

\begin{document}
\maketitle

\noindent\textbf{Status:} OPEN \quad \textbf{No prize}

\noindent\textbf{Tags:} graph theory, chromatic number

\medskip

\noindent\textbf{Problem Statement:}

The cochromatic number of $G$, denoted by $\zeta(G)$, is the minimum number of colours needed to colour the vertices of $G$ such that each colour class induces either a complete graph or empty graph. The dichromatic number of $G$, denoted by $\delta(G)$, is the minimum number $k$ of colours required such that, in any orientation of the edges of $G$, there is a $k$-colouring of the vertices of $G$ such that there are no monochromatic oriented cycles. Must a graph with large chromatic number have large dichromatic number? Must a graph with large cochromatic number contain a graph with large dichromatic number?




The first question is due to Erd\H{o}s and Neumann-Lara. The second question is due to Erd\H{o}s and Gimbel. A positive answer to the second question implies a positive answer to the first via the bound mentioned in [760] .


\bigskip
\noindent\small{Source: \url{https://www.erdosproblems.com/761}}
\end{document}
